近年、Vision-Language-Action Model（VLA）やVision-Language Model（VLM）を用いたロボット操作は、物体を掴む、置く、移動するなどの短期的な作業で大きく進展している。
一方で、机上片付けや家庭内支援のような長期ロボットタスクでは、現状の手法は次の動作を生成することには有効であっても、現在の作業段階で何が不確かであり、その不確かさにどう対処すべきかを適切に判断することが難しい。
本研究の目的は、長期ロボットタスクにおいて、ロボットが現在どの作業段階にあり、それまでにどのような操作を行ってきたかに基づいて、不確かな状態への対処を適切に切り替えるための診断方法を明らかにすることである。
この目的を達成するためには、同じ見え方であっても、現在の作業段階やこれまでの作業履歴に応じて、実行、再観測、人間への確認、安全停止、後回しのいずれを選ぶべきかを判断するための基準を表現する必要がある。
本研究では、ロボットが現在見ている画像、これまでの作業履歴、現在の作業段階、次に取り得る作業候補、各段階で満たすべき条件を入力として、不確かな状態を診断し、その結果に応じて適切な対処行動へ分岐する手法を提案する。
そのために、VLMを次の動作を直接生成するためではなく、判断に必要な情報が十分か、不確かさの原因は何か、その不確かさが作業にどの程度関わるか、どのように対処すべきか、実行時にどのような安全条件が必要かを整理して出力するために用いる。
評価では、同じ視覚情報に対して作業段階だけを変えた対照ペアを設計し、VLMが状況に応じて分岐判断を切り替えられるかを検証する。
また、追加学習を行わないVLMや入力条件を簡略化した場合との比較により、この診断方法の有効性を評価する。
さらに、実機デモを通じて、診断結果が安全な分岐制御に接続可能かを確認する。
本研究により、ロボットが「分からない」状態を一律に扱うのではなく、現在の作業段階やこれまでの作業履歴に応じて、実行する、見直す、聞く、止まる、後回しにする判断を切り替えるための基礎的知見を得ることを目指す。
